% !TEX root = ../../../MathModel.tex
% 负责人：罗财能。
\subsubsection{模型准备}
\label{sec:q1:preparation}

\textbf{（1）数据与计量单位。}
读取原始节点、逐箱需求和运输无人机参数，得到1个调度中心、15个服务区、3类机型及80个不可拆货箱。
货物合计758\,kg、2.011\,m$^3$。
按服务区、物资类别、单箱质量和体积建立等价类别；该压缩只交换物理属性相同的箱子，求解后仍回填原始编号。
组批可行性计算中将质量和体积转换为整数克、立方厘米，避免容量边界的累加误差；飞行距离、时间与能耗分别使用m、s与kWh。
本问不把机队库存、电池周转、送达时限和通信约束加入模型。

\textbf{（2）地形航段。}
以DEM覆盖区中心为原点，采用WGS84椭球在中心纬度的经向、纬向米制比例将经纬度仿射转换为局部坐标。
这种短距离平面近似保持航路与原始栅格的相交关系。
记单程水平距离为$d_i$，调度中心地面海拔为$z_O$，服务区地面海拔为$z_i$，航线穿过或触及的全部DEM像元集合为$\mathcal C_i$。
对每个像元采用原始栅格高程$z_c$，巡航海拔为
\begin{equation}
 H_i=\max\left\{\max_{c\in\mathcal C_i}z_c+50,\ z_O,\ z_i+30\right\}.
 \label{eq:q1:cruise}
\end{equation}
采用栅格边界交点枚举整条航线经过的像元，包括恰好触及边界或角点的像元，避免定距采样漏掉短距离穿过的高地。
式\eqref{eq:q1:cruise}的端点项保证爬升与下降高度非负；本题单点往返线路按沿线最高地形加50\,m确定巡航海拔。
去程与返程爬升量分别为$u_i=H_i-z_O$和$v_i=H_i-z_i-30$，去程下降量为$v_i$，返程下降量为$u_i$。

\textbf{（3）补充物理口径。}
去程携带本架次全部货物，交付后返程为空载。
沿用题设载荷相关等效航程结构，水平能耗按可用电量与航程比例计算，爬升附加能耗按重力势能除以效率计算。
重力加速度统一取$g_0=9.80665$\,m/s$^2$；下降不另计能耗或能量回收，交接阶段不添加题目未提供的悬停功率。
这些补充口径供各问题一致使用，不代表风雨环境下的实测耗电模型。
